
\section*{3. uzdevums}

Sākuma tuvinājumam izvēlamies:
\[
\theta=0{,}5\ \mathrm{rad},
\qquad
A_{\pm1}^{(0)}=\frac{\theta}{2}=0{,}25,
\qquad
A_{n\neq\pm1}^{(0)}=0.
\]
Ievietojot šos koeficientus Furjē summā:
\[
\alpha^{(0)}(\tau)
=0{,}25e^{-i\Omega\tau}+0{,}25e^{i\Omega\tau}
=0{,}5\cos(\Omega\tau),
\qquad
\Omega=\frac{1}{1{,}3}.
\]

\begin{figure}[H]
    \centering
    \includegraphics[width=0.9\linewidth]{art/3uzd.png}
    \caption{Sākuma tuvinājuma $\alpha^{(0)}(\tau)$ reālā
    un imaginārā daļa viena perioda laikā.}
    \label{fig:3uzd}
\end{figure}

Grafikā redzams, ka sākuma tuvinājums ir kosinusa funkcija
ar amplitūdu $0{,}5\ \mathrm{rad}$.
Pie $\tau=0$ svārsts atrodas maksimālajā novirzē:
\[
\alpha^{(0)}(0)=0{,}5,
\qquad
\left.\frac{d\alpha^{(0)}}{d\tau}\right|_{\tau=0}=0.
\]
Tātad sākuma tuvinājums atbilst palaišanai bez sākuma ātruma.
Pie $\tau=\widetilde{T}/2$ novirze ir
$-0{,}5\ \mathrm{rad}$, bet pēc pilna perioda
$\widetilde{T}=2{,}6\pi$ funkcija atgriežas sākuma vērtībā.

Imaginārā daļa ir nulle, jo koeficienti
$A_{-1}^{(0)}=A_1^{(0)}$ ir reāli un eksponentu
imaginārās daļas savstarpēji atceļas.
Šis vēl ir tikai sākuma tuvinājums, nevis nelineārā
kustības vienādojuma gala atrisinājums.

\newpage
\section*{4. uzdevums}

Vispirms aprēķinām sākuma tuvinājuma sinusu:
\[
\sin\bigl(\alpha^{(0)}(\tau)\bigr)
=\sin\bigl(0{,}5\cos(\Omega\tau)\bigr).
\]
Izmantojot Furjē transformāciju, atrodam koeficientus
$B_n^{(0)}$, bet no tiem aprēķinām nākamo tuvinājumu:
\[
A_n^{(1)}=\frac{B_n^{(0)}}{\Omega^2n^2},
\qquad n\neq0.
\]
Pie $n=0$ šo formulu nevar izmantot, jo tā satur dalīšanu
ar nulli. Svārstībām ap līdzsvara stāvokli izvēlamies
$A_0=0$; arī $B_0$ šajā aprēķinā ir nulle skaitliskās
precizitātes robežās.

Aprēķinos izmantojam $M=4096$ vienmērīgi izvietotus
laika punktus vienā periodā. Tabulā parādīti pirmās
iterācijas un gala rezultāta koeficienti.
Pozitīvajiem un negatīvajiem indeksiem koeficienti
sakrīt skaitliskās precizitātes robežās.

\begin{table}[H]
    \centering
    \small
    \begin{tabular}{c|c|c|c}
        $n$ & $B_n^{(0)}$ & $A_n^{(1)}$ & $A_n^{(k)}$ \\
        \hline
$0$ & $0$ & $0$ & $0$ \\
$\pm 1$ & $0{,}2422684577$ & $0{,}4094336935$ & $0{,}9894564073$ \\
$\pm 2$ & $0$ & $0$ & $0$ \\
$\pm 3$ & $-0{,}002563729995$ & $-0{,}0004814115212$ & $-0{,}02471884678$ \\
$\pm 4$ & $0$ & $0$ & $0$ \\
$\pm 5$ & $8{,}053627241\cdot 10^{-6}$ & $5{,}444252015\cdot 10^{-7}$ & $0{,}001039490636$ \\
$\pm 6$ & $0$ & $0$ & $0$ \\
$\pm 7$ & $-1{,}201586733\cdot 10^{-8}$ & $-4{,}144248119\cdot 10^{-10}$ & $-5{,}202157798\cdot 10^{-5}$ \\
$\pm 8$ & $0$ & $0$ & $0$ \\
$\pm 9$ & $1{,}044677144\cdot 10^{-11}$ & $2{,}179635029\cdot 10^{-13}$ & $2{,}833728062\cdot 10^{-6}$ \\
    \end{tabular}
    \caption{Sākuma tuvinājuma sinusa Furjē koeficienti,
    nākamā tuvinājuma koeficienti un koeficienti pēc konverģences.}
\end{table}

Atkārtojam Furjē transformāciju un koeficientu
pārrēķināšanu, līdz:
\[
\max_n\left|A_n^{(k+1)}-A_n^{(k)}\right|<10^{-13}.
\]
Izmantojot $N=9$ jeb indeksus $n=-N,\ldots,N$,
konverģence sasniegta pēc 24 iterācijām.
Pēdējo divu tuvinājumu lielākā koeficientu starpība:
\[
\Delta A_{\max}=4{,}463097819\cdot 10^{-14}.
\]
Šī starpība raksturo iterāciju konverģenci.
Harmoniku skaita ietekmi uz amplitūdas precizitāti
pārbaudām atsevišķi.

Vislielāko ieguldījumu dod pamat­harmonika
$n=\pm1$. Nākamo lielāko ieguldījumu dod
$n=\pm3$, bet augstāko nepāra harmoniku
koeficienti ir arvien mazāki.
Harmonikas $n$ un $-n$ kopā veido vienu
reālu kosinusa locekli:
\[
A_ne^{in\Omega\tau}+A_{-n}e^{-in\Omega\tau}
=2A_n\cos(n\Omega\tau).
\]
Tādēļ iegūto atrisinājumu var pierakstīt:
\[
\begin{aligned}
\alpha(\tau)\approx {} &1{,}978912815\cos(\Omega\tau) \\
&- 0{,}04943769355\cos(3\Omega\tau) \\
&+ 0{,}002078981273\cos(5\Omega\tau) \\
&- 0{,}000104043156\cos(7\Omega\tau) \\
&+ 5{,}667456125\cdot 10^{-6}\cos(9\Omega\tau)
\end{aligned}
\]

Pāra harmonikas $n=\pm2,\pm4,\ldots$ ieguldījumu
nedod; to skaitliski aprēķinātās vērtības ir nulle
noapaļošanas precizitātes robežās.
Tas izriet no svārstību simetrijas:
\[
\alpha\left(\tau+\frac{\widetilde{T}}{2}\right)
=-\alpha(\tau).
\]
Pēc pusperioda novirze maina zīmi, tāpēc Furjē
izvirzījumā paliek tikai nepāra harmonikas.
Arī pastāvīgais loceklis $A_0$ ir nulle.

\begin{figure}[H]
    \centering
    \includegraphics[width=0.9\linewidth]{art/4uzd.png}
    \caption{Aprēķinātā $\alpha(\tau)$ reālā un
    imaginārā daļa pēc koeficientu konverģences.}
    \label{fig:4uzd}
\end{figure}

Grafikā redzamas periodiskas svārstības ap līdzsvara
stāvokli $\alpha=0$.
Atšķirībā no sākuma tuvinājuma gala funkcija
nav tikai viens kosinuss, jo tai ir arī augstāko
nepāra harmoniku ieguldījums.
Imaginārā daļa ir nulle skaitliskās precizitātes
robežās, tāpēc iegūtais atrisinājums ir reāls.

Svārstību amplitūdu nosakām pozitīvajā galējā
stāvoklī pie $\tau=0$:
\[
\theta_N=\alpha(0)=\sum_{n=-N}^{N}A_n
=1{,}931455727\ \mathrm{rad}.
\]
Sākumā izvēlētie $0{,}5\ \mathrm{rad}$ ir tikai
iterāciju sākuma tuvinājuma amplitūda.
Aprēķinā ir dots periods, tāpēc gala amplitūdu
nosaka prasība iegūt svārstības ar
$\widetilde{T}=1{,}3\cdot2\pi$.

Lai pārbaudītu, cik harmoniku vajag sešiem
zīmīgajiem cipariem, atkārtojam aprēķinu pie
dažādām $N$ vērtībām.
Salīdzinājumam izmantojam rezultātu ar $N=40$.
Palielinot $N$ no 30 līdz 40, amplitūdas izmaiņa ir:
\[
|\theta_{40}-\theta_{30}|=0\ \mathrm{rad}.
\]

\begin{table}[H]
    \centering
    \begin{tabular}{c|c|c}
        $N$ & $\theta_N$, rad & $|\theta_N-\theta_{40}|$, rad \\
        \hline
$1$ & $1{,}957384118$ & $0{,}02592869588$ \\
$2$ & $1{,}957384118$ & $0{,}02592869588$ \\
$3$ & $1{,}929434949$ & $0{,}002020472719$ \\
$4$ & $1{,}929434949$ & $0{,}002020472719$ \\
$5$ & $1{,}931551553$ & $9{,}613107654\cdot 10^{-5}$ \\
$6$ & $1{,}931551553$ & $9{,}613107654\cdot 10^{-5}$ \\
$7$ & $1{,}931450105$ & $5{,}316799437\cdot 10^{-6}$ \\
$8$ & $1{,}931450105$ & $5{,}316799437\cdot 10^{-6}$ \\
$9$ & $1{,}931455727$ & $3{,}049804407\cdot 10^{-7}$ \\
$10$ & $1{,}931455727$ & $3{,}049804405\cdot 10^{-7}$ \\
    \end{tabular}
    \caption{Aprēķinātās amplitūdas salīdzinājums pie
    dažāda lielākā harmonikas indeksa $N$.}
\end{table}

Šajā salīdzinājumā pirmais $N$, pie kura amplitūdas
noapaļojums līdz sešiem zīmīgajiem cipariem sakrīt
ar salīdzinājuma rezultātu un absolūtā starpība ir
mazāka par pusi no pēdējā cipara vienības, ir:
\[
N=9,
\qquad
\theta\approx1{,}93146\ \mathrm{rad}.
\]
Tātad šajā skaitliskajā pārbaudē pietiek Furjē
summu paturēt līdz $|n|=9$.
Jāņem vērā, ka pāra harmoniku koeficienti ir nulle,
tāpēc faktisko ieguldījumu dod tikai nepāra harmonikas.
\newpage
